\section{项目出发点与设计目标}
\label{sec:positioning}
\subsection{从一笔指令到一份适合当下的安排}
银行提供的操作能力日益丰富，但用户仍要在收入、账单、家庭责任与生活目标之间，自己拼出一份安排。“帮我打理这两万元”对现金紧张者和储备充足者，不应产生相同的行动；昨天能够执行的计划，也可能因工资延期而在今天失去可行性。TROPOS 的出发点，是让银行智能体先理解资金在生活中的位置，再决定可以准备什么操作。

\begin{keybox}[核心设计]
以动态画像形成安排，以确定性内核守住执行边界。模型提出候选，用户保留决定权，账本事实由服务端核验。
\end{keybox}

“千人千策，随境而变”在工程上对应两个要求。一是个体差异进入决策条件：现金底线、必要支付、期限与用户主动声明的偏好，影响建议空间和核验结果。二是状态变化进入后续决策：收入到账、延期、账单变化或画像更正，会改变账户状态或版本，使旧方案接受重新检查。系统不把一次问卷变成永久标签，也不因为一次收入冲击擅自改变用户的长期目标。

项目面向比赛评审和可复现的技术验证，核心问题是如何把模型生成与金融状态变更衔接起来：允许模型理解自然语言与提出候选安排，同时以确定性业务服务掌握事实与执行权。本文将产品立意落实到数据结构、约束、确认过程和回执，而不是以功能数量代替这一主线。

\subsection{四种资金需要与项目版图}
Protect、Recover、Accumulate、Grow 分别对应现金流守护、负债减压、财富积累与长期资金安排。它们是随处境变化的任务重点，可以并存，也可以调整；不能依据年龄、职业等身份直接为用户分配固定金融策略。

\begin{table}[htbp]\centering\small
\caption{资金需要与当前实现的对应关系}
\begin{tabularx}{\textwidth}{@{}L{2.2cm}YY@{}}\toprule
任务重点 & 要回答的问题 & 当前可体验的支撑 \\\midrule
现金流守护 & 必要账单和生活缓冲能否保留 & 30 天现金路径、现金底线、事件冲击复核 \\
负债减压 & 到期义务怎样排序与核对 & 还款与账单汇总、订阅复核、多步骤任务 \\
财富积累 & 日常之外能为目标留出多少空间 & 目标储蓄草案、生活安排比较、预算核验 \\
长期资金安排 & 用钱期限与风险约束是否匹配 & 主动风险声明、期限与流动性核对、理财沙箱 \\\bottomrule
\end{tabularx}
\end{table}

产品由三层场景连接起来：银行事务提供转账、查询、订阅及卡片等操作工具；现金规划判断近期约束与目标空间；生活联动把家庭、购物、教育、健康和出行等需要放在同一份资金日历中。长期愿景是持续适应人的财务变化；当前落点仍是有限期限、明确授权、可核验的比赛原型。

\subsection{赛题对应与实现边界}
赛事国内人工智能方向聚焦银行 AI 智能体，通过自然语言理解意图、规划并办理业务，兼顾安全与体验。TROPOS 以转账、账单分析与订阅办理为基础，扩展卡片、理财沙箱及生活联动。赛题名称依据赛事官网，访问日期为 2026 年 11 月 2 日\footnote{FinTechathon 赛事官网：\url{https://fintechathon.g-ican.com/}。本文不据官网介绍推定未公开的评分权重或宣称已经通过赛事验收。}。

银行沙箱中的账户、余额、交易、复式分录和银行回执均为合成记录，不连接真实银行清算。美团官方链路、外部模型和 SMTP 则可能产生真实的对外请求；它们分别受连接器、数据边界与部署配置约束。真实美团订单创建不等于合成银行付款，跳转收银台也不等于支付成功。

系统当前没有实名核验、银行卡要素认证、生产清算、持续后台自动付款或经过验证的完整生命周期模型。页面中的资金预测是基于已知输入的条件计算，不是工资到账保证、收益承诺或完整财务风险评估。
